\section{Tensor network decomposition}

We now derive a tensor network decomposition of the tensor $F^{\catorder}[\catvariableof{[\catorder]},\headvariableof{[\catorder]}]$.


\begin{lemma}
    The Fourier transform of basis vectors $\onehotmapofat{\shortcatindices}{\shortcatvariables}$ is the elementary tensor
    \begin{align*}
        \frac{1}{\sqrt{2^{\catorder}}} \bigotimes_{\catenumeratorin} \left( \fbasisat{\headvariableof{\catenumerator}} + \expof{2\pi i \cdot (0.\catindexof{\catorder-\catenumerator-1}\cdots\catindexof{\catorder-1})}\cdot \tbasisat{\headvariableof{\catenumerator}}\right)
    \end{align*}
\end{lemma}
\begin{proof}
    For each $\headindexof{[\catorder]}$ we have
    \begin{align*}
        &\contraction{
            \{ \left( \fbasisat{\headvariableof{\catenumerator}} + \expof{2\pi i \cdot 0.\catindexof{\catorder-\catenumerator-1}\cdots\catindexof{\catorder-1}}\cdot \tbasisat{\headvariableof{\catenumerator}}\right) \wcols \catenumeratorin\}
            \cup \{\onehotmapofat{\headindexof{[\catorder]}}{\headvariableof{[\catorder]}}\}
        } \\
        &\quad= \prod_{\catenumeratorin} \expof{2\pi i \cdot (0.\catindexof{\catorder-\catenumerator-1}\cdots\catindexof{\catorder-1}) \cdot \headindexof{\catenumerator}} \\
        &\quad= \prod_{\catenumeratorin} \expof{2\pi i \cdot (\catindexof{0}\cdots\catindexof{\catorder-\catenumerator-2}.\catindexof{\catorder-\catenumerator-1}\cdots\catindexof{\catorder-1}) \cdot \headindexof{\catenumerator}} \\
        &\quad= \expof{2\pi i \cdot \catindex \sum_{\catenumeratorin} \headindexof{\catenumerator} \cdot 2^{\catorder-\catenumerator-1}} \\
        &\quad= \expof{2\pi i \cdot  \frac{\catindex\cdot\headindex}{2^{\catorder}}} \, .
    \end{align*}
    Since the basis coefficients of the elementary tensor coincide with those of the Fourier transform of $\onehotmapofat{\shortcatindices}{\shortcatvariables}$, both tensors are equal.
\end{proof}


We introduce the diagonal unitary
\begin{align*}
    U^{\catorder}[\indexedshortcatvariables,\headvariableof{[\catorder]}=\headindexof{[\catorder]}]
    = \begin{cases}
          \expof{2\pi i \cdot \left(\sum_{\catenumeratorin} \catindexof{\catenumerator} \cdot 2^{-\catenumerator-2} \right)} & \ifspace \shortcatindices=\headindexof{[\catorder]} \\
          0 & \elsetext
    \end{cases}
\end{align*}
Note that $U^{\catorder}$ is diagonal and the tensor product of matrices with variables $\catvariableof{\catenumerator},\headvariableof{\catenumerator}$.
Diagonality can be used to contract $U^{\catorder}$ to an arbitrary tensor with computational demand in $\mathcal{0}(2^{\catorder})$.

Using diagonality and elementarity, we have
\stantikzpic{
    \draw (0,2) -- (0,3);
    \draw (-4.5,3) rectangle (4.5,6);
    \drawtextnode{0}{4.5}{$\begin{pmatrix}
                               1 \\ \expof{2\pi i \cdot 2^{-\catorder+1}}
    \end{pmatrix}$}
    \drawvariabledot{0}{2}
    \drawhwire{-6}{2}{$\tilde{\headvariable}_{\catorder-1}$}
    \draw (-4,2) -- (5,2);
    \drawsmalltensor{6}{2}{$\identity$}
    \drawhwire{7}{2}{$\headvariableof{\catorder}$}
    \drawtextnode{0}{0.5}{$\vdots$}
    \drawhwire{-6}{-2}{$\tilde{\headvariable}_0$}
    \draw (-4,-2) -- (5,-2);
    \drawsmalltensor{6}{-2}{$\identity$}
    \drawhwire{7}{-2}{$\headvariableof{1}$}
    \drawvariabledot{0}{-2}
    \draw (0,-2) -- (0,-3);
    \draw (-4.5,-3) rectangle (4.5,-6);
    \drawtextnode{0}{-4.5}{$\begin{pmatrix}
                                1 \\ \expof{2\pi i \cdot 2^{-2}}
    \end{pmatrix}$}

    \drawtextnode{-8}{0}{$=$}

    \begin{scope}[shift={(-18,0)}]
        \drawhwire{0}{2}{$\tilde{\headvariable}_{\catorder-1}$}
        \drawtextnode{1}{0.5}{\colorlabelsize $\vdots$}
        \drawhwire{0}{-2}{$\tilde{\headvariable}_{0}$}
        \draw (2,-3) rectangle (4,3);
        \drawtextnode{3}{0}{\corelabelsize $U^{\catorder}$}

        \drawhwire{4}{2}{$\headvariableof{\catorder}$}
        \drawtextnode{5}{0.5}{\colorlabelsize $\vdots$}
        \drawhwire{4}{-2}{$\headvariableof{1}$}

    \end{scope}
}


\begin{lemma}
    The Fourier transform on $[2^{\catorder}]$ can be decomposed into a Fourier transform on $[2^{(\catorder-1)}]$ as:
    \stantikzpic{
        \drawhwire{-3}{2.5}{$\catvariableof{0}$}
        \drawtextnode{-2}{0.5}{\colorlabelsize $\vdots$}
        \drawhwire{-3}{-2.5}{$\catvariableof{\catorder-1}$}
        \draw (-1,-3) rectangle (1,3);
        \drawtextnode{0}{0}{\corelabelsize $F^{\catorder}$}
        \drawhwire{1}{2.5}{$\headvariableof{\catorder-1}$}
        \drawtextnode{2}{0.5}{\colorlabelsize $\vdots$}
        \drawhwire{1}{-2.5}{$\headvariableof{0}$}

        \drawtextnode{8}{0}{$=$}

        \begin{scope}[shift={(15,0)}]
            \drawhwire{-4}{2.5}{$\catvariableof{0}$}
            \drawtextnode{-3}{1.5}{\colorlabelsize $\vdots$}
            \drawhwire{-4}{-0.5}{$\catvariableof{\catorder-2}$}

            \draw (-2,-1) rectangle (0,3);
            \drawtextnode{-1}{1}{\corelabelsize $F^{\catorder-1}$}

            \draw (-4,-2) -- (5,-2);
            \drawsmalltensor{6}{-2}{$\hgate$}
            \drawhwire{7}{-2}{$\headvariableof{0}$}
            \draw (3,-2) -- (3,-1);
            \drawvariabledot{3}{-2}

            \drawhwire{0}{2.5}{$\tilde{\headvariable}_{\catorder-2}$}
            \drawtextnode{1}{1.5}{\colorlabelsize $\vdots$}
            \drawhwire{0}{-0.5}{$\tilde{\headvariable}_{0}$}
            \draw (2,-1) rectangle (4,3);
            \drawtextnode{3}{1}{\corelabelsize $U^{\catorder-1}$}
            \draw (4,2.5) -- (7,2.5);
            \drawtextnode{5}{1.5}{\colorlabelsize $\vdots$}
            \drawhwire{7}{2.5}{$\headvariableof{\catorder-1}$}
            \draw (4,-0.5) -- (7,-0.5);
            \drawhwire{7}{-0.5}{$\headvariableof{1}$}

        \end{scope}

    }
\end{lemma}
\begin{proof}
    It suffices to show that the Fourier transform of arbitrary basis tensors $\onehotmapofat{\shortcatindices}{\shortcatvariables}$ follows the decomposition.
    The Fourier transform of $\onehotmapofat{[\catorder-1]}{\catvariableof{[\catorder-1]}}$ is by the Lemma above
    \begin{align*}
        \frac{1}{\sqrt{2^{\catorder-1}}} \bigotimes_{\catenumerator\in[\catorder-1]}
        \left( \fbasisat{\tilde{\headvariable}_{\catenumerator}} + \expof{2\pi i \cdot (0.\catindexof{\catorder-\catenumerator-2}\cdots\catindexof{\catorder-2})}\cdot \tbasisat{\tilde{\headvariable}_{\catenumerator}}\right)
    \end{align*}

    %Case $\catindexof{\catorder-1}=0$
    In case of $\catindexof{\catorder-1}=0$, the Fourier transform of $\onehotmapofat{\shortcatindices}{\shortcatvariables}$ is the tensor product of this state with $\hplusbasat{\catvariableof{0}}$.
    This is reproduced by the decomposition, since the controlled unitary $U^{\catorder-1}$ acts trivially and the Hadamard gate transforms $\fbasisat{\catorder-1}$ to $\hplusbasat{\catvariableof{0}}$.

    %Case $\catindexof{\catorder-1}=1$
    In case of $\catindexof{\catorder-1}=1$, the leg vectors $\left( \fbasisat{\tilde{\headvariable}_{\catenumerator}} + \expof{2\pi i \cdot (0.\catindexof{\catorder-\catenumerator-2}\cdots\catindexof{\catorder-2})}\cdot \tbasisat{\tilde{\headvariable}_{\catenumerator}}\right)$ are further transformed into
    \begin{align*}
        \left( \fbasisat{\headvariableof{\catenumerator+1}} + \expof{2\pi i \cdot (0.\catindexof{\catorder-\catenumerator-2}\cdots\catindexof{\catorder-2}\catindexof{\catorder-1})}\cdot \tbasisat{\headvariableof{\catenumerator+1}}\right) \, .
    \end{align*}
    In both cases we arrive at the Fourier transform of $\onehotmapofat{\shortcatindices}{\shortcatvariables}$.
\end{proof}

By iteratively decomposing $F^{\catorder-1}$ until we arrive at $F^{1}$ (which is the $\hgate$ matrix), we get a decomposition of $F^{\catorder}$ into $\catorder$ $\hgate$ matrices and controlled $U$ unitaries.

